% ===============================================================
%  جلسه بیست‌وسوم
% ===============================================================
\chapter{تخمین بازه‌ای و بازه اطمینان}

اسلایدهای این جلسه ابتدا تخمین‌گر بیشینه درست‌نمایی را مرور می‌کنند (تعریف گسسته و پیوسته، کانال نویزی، کارت‌های داخل جعبه، گاوسی دوپارامتری و گرایش $\hat\sigma^2_\ML$)؛ این مطالب به طور کامل در جلسه‌های ۲۱ و ۲۲ آمده است. موضوع اصلی و جدید این جلسه \textbf{تخمین بازه‌ای} است: به جای اینکه فقط «یک عدد» به عنوان حدس پارامتر گزارش کنیم، \textbf{یک بازه} گزارش می‌کنیم و می‌گوییم با چه اطمینانی پارامتر واقعی در آن قرار دارد. در این فصل مفهوم بازه اطمینان را از پایه می‌سازیم، آن را برای میانگین توزیع گاوسی با واریانس معلوم به دست می‌آوریم و سپس با معرفی \textbf{توزیع $t$ استیودنت}، حالت واریانس نامعلوم را حل می‌کنیم.

\begin{summarybox}[خلاصه جلسه‌های قبل]
\begin{itemize}
\item \lr{MLE}: $\hat\theta_\ML=\argmax_\theta\prod_{i=1}^nf_X(X_i;\theta)$ (در حالت گسسته به جای چگالی، تابع جرم احتمال).
\item کارت‌های داخل جعبه: $\hat n_\ML=X$ (گرایش‌دار) و $2X-1$ (بی‌گرایش).
\item گاوسی با $\mu$ و $\sigma^2$ نامعلوم:
\[
\hat\mu_\ML=\bar X_n,\qquad \hat\sigma^2_\ML=\frac1n\sum_{i=1}^n\big(X_i-\bar X_n\big)^2,\qquad \E\hat\sigma^2_\ML=\frac{n-1}{n}\sigma^2 .
\]
\item تخمین‌گر بی‌گرایش واریانس (واریانس نمونه‌ای):
\[
\bar S_n^2\triangleq\frac{1}{n-1}\sum_{i=1}^n\big(X_i-\bar X_n\big)^2,\qquad \E\bar S^2_n=\sigma^2 .
\]
\end{itemize}
\end{summarybox}

\section{چرا تخمین نقطه‌ای کافی نیست؟}

فرض کنید با $n=25$ اندازه‌گیری، میانگین نمونه‌ای را $\bar x=52.3$ به دست آورده‌ایم و آن را به عنوان تخمین میانگین واقعی $\mu$ گزارش می‌کنیم. دو پرسش طبیعی پیش می‌آید:
\begin{itemize}
\item \textbf{این عدد چقدر قابل اعتماد است؟} اگر آزمایش را تکرار کنیم، $\bar x$ احتمالاً عدد دیگری مثل $50.8$ یا $53.9$ خواهد شد. خود عدد $52.3$ هیچ اطلاعاتی درباره این «لرزش» نمی‌دهد.
\item \textbf{احتمال درست بودن آن چقدر است؟} اگر $\bar X_n$ متغیری پیوسته باشد، احتمال اینکه \textbf{دقیقاً} برابر $\mu$ شود صفر است: $\Prob(\bar X_n=\mu)=0$. پس تخمین نقطه‌ای «تقریباً همیشه» غلط است؛ سؤال درست این است که \textbf{چقدر} غلط است.
\end{itemize}
در زندگی روزمره هم همین را می‌بینیم: نتیجه یک نظرسنجی را به شکل «۴۵٪ با خطای ±۳٪» گزارش می‌کنند، نه فقط «۴۵٪». تخمین بازه‌ای دقیقاً همین ایده را صورت‌بندی می‌کند: \textbf{یک بازه} به همراه \textbf{میزان اطمینان} به آن.

\section{صورت‌بندی تخمین بازه‌ای}

\begin{definition}[مسئله تخمین بازه‌ای \lr{(Interval Estimation)}]
فرض کنید از توزیع پارامتری $f_X(x;\theta)$ (که $\theta$ در آن \textbf{غیرتصادفی ولی نامعلوم} است) نمونه‌های \lr{i.i.d.} $X_1,\ldots,X_n$ در دست است. برای یک عدد دلخواه $\delta\in[0,1]$ می‌خواهیم دو تابع از نمونه‌ها
\[
\theta_l=\theta_l\big(X_{1:n}\big),\qquad \theta_u=\theta_u\big(X_{1:n}\big)
\]
بیابیم به طوری که
\[
\Prob\Big(\theta_l\big(X_{1:n}\big)\le\theta\le\theta_u\big(X_{1:n}\big)\Big)\ \ge\ 1-\delta ,
\]
و در عین حال طول بازه، یعنی $\theta_u-\theta_l$، تا حد امکان \textbf{کوچک} باشد. مقدار $\delta$ بسته به کاربرد انتخاب می‌شود (مقادیر رایج: $0.1$، $0.05$ و $0.01$).
\end{definition}

این تعریف دو خواسته متضاد را کنار هم گذاشته است:
\begin{itemize}
\item \textbf{اطمینان بالا:} بازه باید با احتمال زیاد پارامتر را «بپوشاند».
\item \textbf{دقت بالا:} بازه باید کوتاه باشد.
\end{itemize}
اگر فقط خواسته اول را داشتیم، جواب بدیهی $(\theta_l,\theta_u)=(-\infty,+\infty)$ با احتمال پوشش $1$ همیشه کار می‌کرد، ولی هیچ اطلاعاتی نمی‌داد! اگر فقط خواسته دوم را داشتیم، بازه به طول صفر (همان تخمین نقطه‌ای) را انتخاب می‌کردیم که احتمال پوشش آن (در حالت پیوسته) صفر است. تخمین بازه‌ای یعنی \textbf{کوتاه‌ترین بازه‌ای که اطمینان خواسته‌شده را تضمین کند}.

\begin{warning}
\warntitle \textbf{چه چیزی تصادفی است؟} در عبارت $\Prob(\theta_l\le\theta\le\theta_u)$، پارامتر $\theta$ یک عدد ثابت (هرچند نامعلوم) است و \textbf{تصادفی نیست}. آنچه تصادفی است \textbf{دو سر بازه} است، چون توابعی از نمونه‌های تصادفی‌اند. به زبان فضای نمونه:
\[
\Prob\Big(\Big\{\omega:\ \theta_l\big(X_{1:n}(\omega)\big)\le\theta\le\theta_u\big(X_{1:n}(\omega)\big)\Big\}\Big)\ \ge\ 1-\delta .
\]
یعنی «احتمال اینکه \textbf{بازه تصادفی} پارامتر ثابت را بپوشاند». (در یکی از اسلایدها در عبارت بالا هر دو طرف به اشتباه $\theta_l$ نوشته شده است؛ طرف راست باید $\theta_u$ باشد.)
\end{warning}

\subsection{تفسیر فراوانی: پوشش در تکرار آزمایش}
یک راه شهودی برای فهم این احتمال: تصور کنید آزمایش (گرفتن $n$ نمونه و ساختن بازه) را بارها و بارها تکرار کنیم. در هر تکرار، بازه جدیدی به دست می‌آید که جایش کمی جابه‌جا شده است. اگر سطح اطمینان $1-\delta=0.95$ باشد، در درازمدت حدود ۹۵ درصد این بازه‌ها مقدار واقعی $\theta$ را در بر می‌گیرند و حدود ۵ درصد آن را از دست می‌دهند. شکل زیر همین را نشان می‌دهد.

\begin{center}
\begin{tikzpicture}
\end{tikzpicture}\\
{\small ۴۰ بار تکرار یک آزمایش: در هر تکرار $n=25$ نمونه از $\mathcal N(50,10^2)$ گرفته شده و بازه اطمینان ۹۵٪ (با $\sigma$ معلوم) ساخته شده است. خط عمودی مقدار واقعی $\mu=50$ است. بازه‌هایی که $\mu$ را نمی‌پوشانند با رنگ قرمز مشخص شده‌اند؛ حدود ۵٪ بازه‌ها چنین‌اند. توجه کنید که $\mu$ ثابت است و این بازه‌ها هستند که تصادفی جابه‌جا می‌شوند.}
\end{center}

\begin{warning}
\warntitle پس از اینکه داده‌ها مشاهده شدند و مثلاً بازه $[48.4,\ 56.2]$ به دست آمد، دیگر هیچ چیز تصادفی‌ای باقی نمانده است: $\mu$ یا در این بازه هست یا نیست (احتمال یا ۱ است یا ۰). عبارت «با احتمال ۹۵٪ مقدار $\mu$ در بازه $[48.4,\ 56.2]$ است» به معنای دقیق کلمه نادرست است. جمله درست این است: «این بازه با روشی ساخته شده که در ۹۵٪ تکرارها پارامتر را می‌پوشاند.» (در تخمین بیزی، جلسه ۲۴، که خود پارامتر تصادفی فرض می‌شود، جمله اول معنا پیدا می‌کند.)
\end{warning}

\subsection{تخمین نقطه‌ای در برابر تخمین بازه‌ای}
از نظر مفهومی تفاوت بنیادینی میان تخمین نقطه‌ای و بازه‌ای وجود ندارد:
\begin{itemize}
\item در تخمین نقطه‌ای، اگر $\mathrm{Bias}$ و $\Var$ تخمین‌گر $\hat\theta$ هر دو کوچک باشند، آنگاه $\hat\theta$ «با احتمال زیاد» نزدیک مقدار واقعی است.
\item در تخمین بازه‌ای، همین جمله را \textbf{کمّی} می‌کنیم: به ازای یک $\delta$ ثابت، کران‌هایی پیدا می‌کنیم که «با احتمال حداقل $1-\delta$» پارامتر را در بر بگیرند.
\end{itemize}
در واقع از هر تخمین‌گر نقطه‌ای با خطای میانگین مربعات معلوم می‌توان یک بازه ساخت. طبق نامساوی مارکوف برای متغیر نامنفی $(\hat\theta-\theta)^2$ (جلسه ۱۸):
\[
\Prob\big(|\hat\theta-\theta|\ge\varepsilon\big)=\Prob\big((\hat\theta-\theta)^2\ge\varepsilon^2\big)\le\frac{\E(\hat\theta-\theta)^2}{\varepsilon^2}=\frac{\mathrm{MSE}(\hat\theta)}{\varepsilon^2}.
\]
اگر سمت راست را برابر $\delta$ قرار دهیم، یعنی $\varepsilon=\sqrt{\mathrm{MSE}/\delta}$، نتیجه می‌گیریم که بازه $\hat\theta\pm\sqrt{\mathrm{MSE}/\delta}$ با احتمال حداقل $1-\delta$ پارامتر را می‌پوشاند. پس \textbf{تخمین نقطه‌ای خوب (با $\mathrm{MSE}$ کوچک) خودبه‌خود بازه کوتاهی می‌دهد.} این ایده را در جلسه ۲۴ با جزئیات کامل دنبال می‌کنیم.

\subsection{بازه اطمینان و سطح اطمینان}
\begin{definition}[بازه اطمینان \lr{(Confidence Interval)} و سطح اطمینان \lr{(Confidence Level)}]
مجموعه همه بازه‌های بسته روی خط حقیقی را با
\[
\Delta\triangleq\big\{(x,y)\ \big|\ x,y\in\R,\ x\le y\big\}
\]
نشان می‌دهیم. هر تابع
\[
\big(\theta_l,\theta_u\big):\ \R^n\to\Delta
\]
که به هر مجموعه داده $X_{1:n}$ یک بازه نسبت می‌دهد، یک \textbf{تخمین‌گر بازه‌ای} نامیده می‌شود. بازه‌ای که از اعمال این تابع بر داده‌های \textbf{مشاهده‌شده} به دست می‌آید، \textbf{بازه اطمینان} و مقدار $1-\delta$ \textbf{سطح اطمینان} نام دارد.
\end{definition}
مثلاً «بازه اطمینان ۹۵ درصدی» یعنی بازه‌ای که با سطح اطمینان $1-\delta=0.95$ (یعنی $\delta=0.05$) ساخته شده است.

\section{میانگین توزیع گاوسی با واریانس معلوم}

\subsection{صورت مسئله و انتخاب شکل بازه}
فرض کنید $X_1,\ldots,X_n\sim\mathcal N(\mu,\sigma^2)$ به صورت \lr{i.i.d.} باشند، که در آن $\mu$ نامعلوم و $\sigma^2$ \textbf{معلوم} است. می‌خواهیم بازه اطمینانی با سطح $1-\delta$ برای $\mu$ بسازیم.

\textbf{مرکز بازه کجا باشد؟} طبیعی است بازه را حول بهترین تخمین نقطه‌ای که داریم بسازیم، یعنی میانگین نمونه‌ای $\bar X_n$؛ چون:
\begin{itemize}
\item $\bar X_n$ بی‌گرایش است: $\E\bar X_n=\mu$؛
\item $\bar X_n$ همان $\hat\mu_\ML$ است و در میان تخمین‌گرهای بی‌گرایش کمترین واریانس را دارد (پس بازه حول آن کوتاه‌ترین می‌شود)؛
\item توزیع $\bar X_n$ (چنان‌که در ادامه می‌بینیم) گاوسی و \textbf{متقارن} حول $\mu$ است.
\end{itemize}
به دلیل تقارن، بازه را هم \textbf{متقارن} می‌گیریم. چون انحراف معیار $\bar X_n$ برابر $\sigma/\sqrt n$ است، طبیعی است نصف طول بازه را مضربی از همین مقدار بگیریم:
\[
\Big[\ \bar X_n-k\frac{\sigma}{\sqrt n}\ ,\ \bar X_n+k\frac{\sigma}{\sqrt n}\ \Big],
\]
که در آن $k>0$ عددی است که باید طوری تعیین شود که سطح اطمینان دقیقاً $1-\delta$ شود. (در تمرین‌های آخر فصل ثابت می‌کنیم که بازه متقارن، کوتاه‌ترین بازه ممکن با این سطح اطمینان است.)

\subsection{گام ۱: توزیع میانگین نمونه‌ای}
$\bar X_n=\frac1n\sum_{i=1}^nX_i$ ترکیب خطی از متغیرهای گاوسی مستقل است و می‌دانیم (جلسه ۱۶) که ترکیب خطی متغیرهای گاوسی مستقل، گاوسی است. میانگین و واریانس آن:
\[
\E\bar X_n=\frac1n\sum_{i=1}^n\E X_i=\frac1n\cdot n\mu=\mu,\qquad
\Var\bar X_n=\frac1{n^2}\sum_{i=1}^n\Var X_i=\frac{1}{n^2}\cdot n\sigma^2=\frac{\sigma^2}{n}
\]
(در واریانس از استقلال استفاده کردیم: واریانس مجموع متغیرهای مستقل برابر مجموع واریانس‌هاست.) پس
\[
\bar X_n\sim\mathcal N\Big(\mu,\ \frac{\sigma^2}{n}\Big)
\qquad\Longleftrightarrow\qquad
Z\triangleq\frac{\bar X_n-\mu}{\sigma/\sqrt n}\sim\mathcal N(0,1).
\]

\subsection{گام ۲: احتمال اینکه $\bar X_n$ نزدیک $\mu$ باشد}
ابتدا احتمال اینکه $\bar X_n$ در فاصله $k\sigma/\sqrt n$ از $\mu$ بیفتد را حساب می‌کنیم:
\[
\begin{aligned}
\Prob\Big(\mu-k\frac{\sigma}{\sqrt n}\le\bar X_n\le\mu+k\frac{\sigma}{\sqrt n}\Big)
&=F_{\mathcal N(\mu,\sigma^2/n)}\Big(\mu+k\frac{\sigma}{\sqrt n}\Big)-F_{\mathcal N(\mu,\sigma^2/n)}\Big(\mu-k\frac{\sigma}{\sqrt n}\Big)\\
&=F_{\mathcal N(0,1)}(k)-F_{\mathcal N(0,1)}(-k)=\Phi(k)-\Phi(-k)=2\Phi(k)-1 .
\end{aligned}
\]
تساوی دوم با استانداردسازی است: نقطه $\mu+k\sigma/\sqrt n$ پس از کم کردن میانگین $\mu$ و تقسیم بر انحراف معیار $\sigma/\sqrt n$ به نقطه $k$ تبدیل می‌شود. تساوی آخر از تقارن چگالی گاوسی استاندارد است: $\Phi(-k)=1-\Phi(k)$.

\subsection{گام ۳: وارونه کردن رویداد}
نکته کلیدی این است که رویداد بالا را می‌توان طوری بازنویسی کرد که $\mu$ در وسط قرار بگیرد. نامساوی دوطرفه را جدا جدا بررسی می‌کنیم:
\[
\mu-k\frac{\sigma}{\sqrt n}\le\bar X_n\ \Longleftrightarrow\ \mu\le\bar X_n+k\frac{\sigma}{\sqrt n},
\qquad
\bar X_n\le\mu+k\frac{\sigma}{\sqrt n}\ \Longleftrightarrow\ \bar X_n-k\frac{\sigma}{\sqrt n}\le\mu .
\]
پس این دو رویداد \textbf{دقیقاً یکی‌اند} (یک مجموعه از $\omega$ها را توصیف می‌کنند) و احتمال برابر دارند:
\[
\boxed{\ \Prob\Big(\bar X_n-k\frac{\sigma}{\sqrt n}\le\mu\le\bar X_n+k\frac{\sigma}{\sqrt n}\Big)=2\Phi(k)-1\ }
\]
شهود: «$\bar X_n$ در فاصله $d$ از $\mu$ است» با «$\mu$ در فاصله $d$ از $\bar X_n$ است» یکی است؛ فاصله متقارن است!

\subsection{گام ۴: تعیین $k$}
برای اینکه سطح اطمینان دقیقاً $1-\delta$ باشد:
\[
2\Phi(k)-1=1-\delta\quad\Longleftrightarrow\quad\Phi(k)=1-\frac\delta2\quad\Longleftrightarrow\quad k=z_{\delta/2}\triangleq\Phi^{-1}\Big(1-\frac\delta2\Big).
\]
یعنی $k$ نقطه‌ای است که سطح زیر چگالی گاوسی استاندارد در سمت راست آن برابر $\delta/2$ است (و به دلیل تقارن، سمت چپ $-k$ هم $\delta/2$). \textbf{نکته مهم:} معادله تعیین $k$ به $n$، $\mu$ و $\sigma$ وابسته نیست و $k=k(\delta)$ فقط به سطح اطمینان بستگی دارد. علتش این است که کمیت $Z=\frac{\bar X_n-\mu}{\sigma/\sqrt n}$ به ازای همه مقادیر پارامترها توزیع ثابت $\mathcal N(0,1)$ دارد.

\begin{center}
\begin{tikzpicture}
\end{tikzpicture}\\
{\small چگالی گاوسی استاندارد: ناحیه میانی $[-k,k]$ احتمال $1-\delta$ و هر دم احتمال $\delta/2$ دارد. برای $1-\delta=0.95$ داریم $k=1.96$.}
\end{center}

\begin{center}
\renewcommand{\arraystretch}{1.4}
\begin{tabular}{c|c|c|c|c}
سطح اطمینان $1-\delta$ & $0.90$ & $0.95$ & $0.99$ & $0.999$\\\hline
$k=z_{\delta/2}$ & $1.645$ & $1.960$ & $2.576$ & $3.291$
\end{tabular}
\end{center}

\begin{theorem}[بازه اطمینان میانگین گاوسی با واریانس معلوم]
اگر $X_1,\ldots,X_n\sim\mathcal N(\mu,\sigma^2)$ به صورت \lr{i.i.d.} و $\sigma$ معلوم باشد، آنگاه
\[
\Big[\ \bar X_n-z_{\delta/2}\frac{\sigma}{\sqrt n}\ ,\ \bar X_n+z_{\delta/2}\frac{\sigma}{\sqrt n}\ \Big]
\]
یک بازه اطمینان با سطح \textbf{دقیقاً} $1-\delta$ برای $\mu$ است. طول آن $2z_{\delta/2}\,\sigma/\sqrt n$ است.
\end{theorem}

\begin{note}
\notetitle \textbf{رفتار طول بازه.} طول $L=2z_{\delta/2}\sigma/\sqrt n$:
\begin{itemize}
\item با افزایش $n$ مثل $\frac1{\sqrt n}$ کم می‌شود و به صفر میل می‌کند (سازگاری). برای نصف کردن طول بازه، باید تعداد نمونه را \textbf{چهار برابر} کرد.
\item با افزایش $\sigma$ (داده نویزی‌تر) خطی زیاد می‌شود.
\item با افزایش سطح اطمینان (کم شدن $\delta$) زیاد می‌شود: اطمینان بیشتر به قیمت دقت کمتر.
\end{itemize}
\end{note}

\begin{note}
\notetitle \textbf{کمیت محوری \lr{(Pivot)}.} ترفندی که استفاده کردیم یک روش کلی است: کمیتی بسازید که هم به داده و هم به پارامتر وابسته باشد، اما \textbf{توزیعش به پارامتر وابسته نباشد} (اینجا $Z=\frac{\bar X_n-\mu}{\sigma/\sqrt n}\sim\mathcal N(0,1)$). سپس بازه‌ای با احتمال $1-\delta$ برای این کمیت پیدا کنید و نامساوی را برحسب پارامتر وارونه کنید. چنین کمیتی را \textbf{محوری} می‌نامند. در ادامه همین روش را با کمیت محوری $T$ و در تمرین‌ها برای توزیع یکنواخت به کار می‌بریم.
\end{note}

\begin{example}[بازه اطمینان ۹۵٪ با واریانس معلوم]
وزن بسته‌های تولیدی یک کارخانه توزیع گاوسی با انحراف معیار معلوم $\sigma=10$ گرم دارد. وزن $n=25$ بسته اندازه‌گیری شده و میانگین $\bar x=52.3$ گرم به دست آمده است. بازه‌های اطمینان ۹۵٪ و ۹۹٪ را برای میانگین واقعی $\mu$ بیابید.
\begin{solution}
انحراف معیار میانگین نمونه‌ای: $\sigma/\sqrt n=10/\sqrt{25}=10/5=2$.

\textbf{سطح ۹۵٪:} $k=z_{0.025}=1.96$، پس نصف طول بازه $1.96\times2=3.92$ است:
\[
\big[52.3-3.92,\ 52.3+3.92\big]=\big[48.38,\ 56.22\big].
\]
\textbf{سطح ۹۹٪:} $k=z_{0.005}=2.576$، پس نصف طول $2.576\times2=5.152$:
\[
\big[52.3-5.152,\ 52.3+5.152\big]=\big[47.15,\ 57.45\big].
\]
همان‌طور که انتظار داشتیم، اطمینان بیشتر بازه پهن‌تری داده است.
\end{solution}
\end{example}

\begin{example}[تعیین تعداد نمونه لازم]
در مثال قبل، حداقل چند بسته باید وزن شود تا بازه اطمینان ۹۵٪ نصف طولی حداکثر ۱ گرم داشته باشد؟
\begin{solution}
باید $z_{0.025}\,\sigma/\sqrt n\le1$ باشد:
\[
\frac{1.96\times10}{\sqrt n}\le1\quad\Longleftrightarrow\quad\sqrt n\ge19.6\quad\Longleftrightarrow\quad n\ge19.6^2=384.16 .
\]
چون $n$ باید صحیح باشد، $n=385$. مقایسه کنید: با $n=25$ نصف طول $3.92$ بود؛ برای حدود ۴ برابر دقیق‌تر شدن، حدود $16$ برابر نمونه لازم شد ($4^2=16$).
\end{solution}
\end{example}

\begin{note}
\notetitle \textbf{اگر داده‌ها گاوسی نباشند چه؟} اگر $X_i$ها توزیع دلخواهی با واریانس معلوم $\sigma^2$ داشته باشند و $n$ بزرگ باشد، طبق \textbf{قضیه حد مرکزی} (جلسه ۱۷) توزیع $Z=\frac{\bar X_n-\mu}{\sigma/\sqrt n}$ تقریباً $\mathcal N(0,1)$ است؛ پس همان بازه با سطح اطمینان \textbf{تقریبی} $1-\delta$ معتبر است. برای $n$ کوچک و توزیع غیرگاوسی، روش مطمئن‌تر استفاده از نامساوی چبی‌شف است (جلسه ۲۴).
\end{note}

\section{میانگین توزیع گاوسی با واریانس نامعلوم: توزیع $t$ استیودنت}

\subsection{مشکل: $\sigma$ را نمی‌دانیم}
در عمل به ندرت $\sigma$ را می‌دانیم. بازه $\bar X_n\pm z_{\delta/2}\sigma/\sqrt n$ به $\sigma$ نیاز دارد و بدون آن قابل محاسبه نیست. ایده طبیعی این است که $\sigma$ را هم از روی داده تخمین بزنیم و به جای آن، \textbf{انحراف معیار نمونه‌ای} را قرار دهیم:
\[
\bar S_n\triangleq\sqrt{\bar S_n^2},\qquad \bar S_n^2=\frac{1}{n-1}\sum_{i=1}^n\big(X_i-\bar X_n\big)^2 .
\]
اما با این جایگزینی، کمیت
\[
\bar T_n\triangleq\frac{\bar X_n-\mu}{\bar S_n/\sqrt n}
\]
دیگر توزیع $\mathcal N(0,1)$ ندارد! علت: در $Z$ مخرج یک \textbf{عدد ثابت} ($\sigma$) بود، اما در $\bar T_n$ مخرج خودش \textbf{تصادفی} است. گاهی $\bar S_n$ به طور اتفاقی از $\sigma$ کوچک‌تر می‌شود و نسبت را بزرگ می‌کند؛ نتیجه این است که توزیع $\bar T_n$ \textbf{دم‌های سنگین‌تری} از گاوسی دارد (مقادیر بزرگ محتمل‌ترند). این اثر برای $n$ کوچک شدید است و برای $n$ بزرگ (که $\bar S_n$ تقریباً همان $\sigma$ می‌شود) از بین می‌رود.

\subsection{توزیع $t$ استیودنت}
این توزیع را ویلیام سیلی گاست در سال ۱۹۰۸ با نام مستعار «\lr{Student}» منتشر کرد (او در کارخانه آبجوسازی گینس کار می‌کرد و با نمونه‌های کوچک سروکار داشت).

\begin{definition}[توزیع $t$ با $\nu$ درجه آزادی]
اگر $Z\sim\mathcal N(0,1)$ و $V\sim\chi^2_\nu$ (مجموع مربعات $\nu$ متغیر گاوسی استاندارد مستقل) مستقل از هم باشند، توزیع
\[
T=\frac{Z}{\sqrt{V/\nu}}
\]
را \textbf{توزیع $t$ استیودنت با $\nu$ درجه آزادی} می‌نامیم و می‌نویسیم $T\sim t_\nu$. چگالی آن:
\[
f_{t_\nu}(t)=\frac{\Gamma\big(\frac{\nu+1}{2}\big)}{\sqrt{\nu\pi}\ \Gamma\big(\frac\nu2\big)}\Big(1+\frac{t^2}{\nu}\Big)^{-\frac{\nu+1}{2}},\qquad t\in\R .
\]
\end{definition}

شهود تعریف: صورت یک گاوسی استاندارد است و مخرج $\sqrt{V/\nu}$ «تخمینی تصادفی از عدد ۱» است (چون $\E V/\nu=1$). وقتی $\nu$ بزرگ است، طبق قانون اعداد بزرگ $V/\nu\to1$ و $T$ به $Z$ نزدیک می‌شود.

\textbf{ویژگی‌ها:}
\begin{itemize}
\item \textbf{متقارن} حول صفر است (چون $Z$ متقارن است و مخرج مثبت است): $F_{t_\nu}(-k)=1-F_{t_\nu}(k)$.
\item دم‌ها به صورت توانی $|t|^{-(\nu+1)}$ کاهش می‌یابند، نه نمایی مثل گاوسی؛ یعنی \textbf{دم سنگین‌تر}.
\item برای $\nu>2$ واریانس آن $\frac{\nu}{\nu-2}>1$ است (بزرگ‌تر از گاوسی استاندارد).
\item برای $\nu=1$ همان توزیع کوشی است که حتی میانگین ندارد.
\item وقتی $\nu\to\infty$، توزیع $t_\nu$ به $\mathcal N(0,1)$ میل می‌کند.
\end{itemize}

\begin{center}
\begin{tikzpicture}
\end{tikzpicture}\\
{\small چگالی توزیع $t$ با $\nu=1,3,10$ درجه آزادی در مقایسه با گاوسی استاندارد. هرچه $\nu$ کوچک‌تر، قله پایین‌تر و دم‌ها سنگین‌تر.}
\end{center}

\subsection{چرا $\bar T_n$ توزیع $t_{n-1}$ دارد؟}
\begin{theorem}[قضیه نمونه‌گیری از توزیع گاوسی]
اگر $X_1,\ldots,X_n\sim\mathcal N(\mu,\sigma^2)$ به صورت \lr{i.i.d.} باشند، آنگاه:
\begin{enumerate}
\item $\bar X_n\sim\mathcal N(\mu,\sigma^2/n)$؛
\item $\dfrac{(n-1)\bar S_n^2}{\sigma^2}\sim\chi^2_{n-1}$؛
\item $\bar X_n$ و $\bar S^2_n$ \textbf{مستقل} از یکدیگرند.
\end{enumerate}
\end{theorem}
اثبات کامل بندهای ۲ و ۳ از حوصله این درس خارج است، اما شهود آن روشن است:
\begin{itemize}
\item $(n-1)\bar S_n^2/\sigma^2=\sum_i\big(\frac{X_i-\bar X_n}{\sigma}\big)^2$ مجموع $n$ مربع است، ولی این $n$ انحراف آزاد نیستند: همواره $\sum_i(X_i-\bar X_n)=0$. یک قید خطی، یک «درجه آزادی» را می‌گیرد و $n-1$ درجه آزادی باقی می‌ماند. (همین قید دلیل تقسیم بر $n-1$ در $\bar S_n^2$ بود.)
\item استقلال $\bar X_n$ و $\bar S_n^2$ ویژگی خاص توزیع گاوسی است: بردار انحراف‌ها $(X_i-\bar X_n)_i$ با $\bar X_n$ ناهمبسته است، چون
\[
\Cov\big(X_i-\bar X_n,\ \bar X_n\big)=\Cov(X_i,\bar X_n)-\Var\bar X_n=\frac{\sigma^2}{n}-\frac{\sigma^2}{n}=0,
\]
و برای متغیرهای \textbf{توأماً گاوسی}، ناهمبستگی معادل استقلال است (جلسه ۱۴).
\end{itemize}

حالا $\bar T_n$ را طوری بازنویسی می‌کنیم که شکل تعریف $t$ ظاهر شود. صورت و مخرج را بر $\sigma/\sqrt n$ تقسیم می‌کنیم:
\[
\bar T_n=\frac{\bar X_n-\mu}{\bar S_n/\sqrt n}
=\frac{\dfrac{\bar X_n-\mu}{\sigma/\sqrt n}}{\dfrac{\bar S_n}{\sigma}}
=\frac{Z}{\sqrt{\dfrac{(n-1)\bar S_n^2/\sigma^2}{n-1}}}
=\frac{Z}{\sqrt{V/(n-1)}},
\]
که $Z\sim\mathcal N(0,1)$ (بند ۱)، $V=(n-1)\bar S_n^2/\sigma^2\sim\chi^2_{n-1}$ (بند ۲) و این دو مستقل‌اند (بند ۳). پس طبق تعریف:
\[
\boxed{\ \bar T_n=\frac{\bar X_n-\mu}{\bar S_n/\sqrt n}\sim t_{n-1}\ }
\]
نکته بسیار مهم: $\sigma$ از عبارت حذف شد! توزیع $\bar T_n$ به هیچ پارامتر نامعلومی وابسته نیست؛ پس $\bar T_n$ یک \textbf{کمیت محوری} است.

\subsection{ساخت بازه اطمینان}
دقیقاً مثل حالت قبل عمل می‌کنیم. برای $k>0$ و با استفاده از تقارن توزیع $t$:
\[
\Prob\big(-k\le\bar T_n\le k\big)=F_{t,n-1}(k)-F_{t,n-1}(-k)=2F_{t,n-1}(k)-1 .
\]
رویداد $-k\le\frac{\bar X_n-\mu}{\bar S_n/\sqrt n}\le k$ را (با ضرب در عدد مثبت $\bar S_n/\sqrt n$ و جابه‌جایی) وارونه می‌کنیم:
\[
-k\frac{\bar S_n}{\sqrt n}\le\bar X_n-\mu\le k\frac{\bar S_n}{\sqrt n}
\quad\Longleftrightarrow\quad
\bar X_n-k\frac{\bar S_n}{\sqrt n}\le\mu\le\bar X_n+k\frac{\bar S_n}{\sqrt n}.
\]
پس با انتخاب $k$ از معادله $2F_{t,n-1}(k)-1=1-\delta$، یعنی
\[
k=t_{n-1,\delta/2}\triangleq F_{t,n-1}^{-1}\Big(1-\frac\delta2\Big),
\]
داریم:
\[
\boxed{\ \Prob\Big(\bar X_n-t_{n-1,\delta/2}\frac{\bar S_n}{\sqrt n}\le\mu\le\bar X_n+t_{n-1,\delta/2}\frac{\bar S_n}{\sqrt n}\Big)=1-\delta\ }
\]
این بار $k=k(n,\delta)$ به $n$ هم وابسته است (از طریق درجه آزادی)، برخلاف حالت واریانس معلوم. (در یکی از اسلایدها عنوان این بخش به اشتباه «واریانس معلوم» نوشته شده؛ این بخش مربوط به \textbf{واریانس نامعلوم} است.)

\begin{center}
\renewcommand{\arraystretch}{1.4}
\begin{tabular}{c|c|c|c|c|c|c|c|c}
درجه آزادی $\nu=n-1$ & $1$ & $2$ & $4$ & $9$ & $19$ & $29$ & $99$ & $\infty$\\\hline
$t_{\nu,0.025}$ (سطح ۹۵٪) & $12.706$ & $4.303$ & $2.776$ & $2.262$ & $2.093$ & $2.045$ & $1.984$ & $1.960$\\
$t_{\nu,0.005}$ (سطح ۹۹٪) & $63.657$ & $9.925$ & $4.604$ & $3.250$ & $2.861$ & $2.756$ & $2.626$ & $2.576$
\end{tabular}
\end{center}

جدول نشان می‌دهد که برای $n$ کوچک، $k$ بسیار بزرگ‌تر از مقدار گاوسی است (مثلاً با $n=2$ نمونه، $12.706$ به جای $1.96$): این \textbf{بهای ندانستن $\sigma$} است. برای $n\ge30$ تفاوت ناچیز می‌شود.

\begin{example}[بازه اطمینان با واریانس نامعلوم]
طول یک قطعه $n=10$ بار با یک ابزار (با خطای گاوسی) اندازه‌گیری شده است (برحسب میلی‌متر):
\[
9.8,\ 10.2,\ 10.4,\ 9.9,\ 10.0,\ 10.3,\ 9.7,\ 10.1,\ 10.6,\ 9.5 .
\]
بازه اطمینان ۹۵٪ برای طول واقعی $\mu$ را بیابید.
\begin{solution}
\textbf{گام ۱ (میانگین):} مجموع داده‌ها $100.5$ است، پس $\bar x=100.5/10=10.05$.

\textbf{گام ۲ (انحراف‌ها و مربع آن‌ها):}
\[
x_i-\bar x:\ -0.25,\ 0.15,\ 0.35,\ -0.15,\ -0.05,\ 0.25,\ -0.35,\ 0.05,\ 0.55,\ -0.55
\]
(مجموع این انحراف‌ها صفر است، که همان قید خطی است.) مجموع مربعات:
\[
\sum_i(x_i-\bar x)^2=0.0625+0.0225+0.1225+0.0225+0.0025+0.0625+0.1225+0.0025+0.3025+0.3025=1.025 .
\]
\textbf{گام ۳ (انحراف معیار نمونه‌ای):}
\[
\bar s^2=\frac{1.025}{10-1}=0.1139,\qquad \bar s=\sqrt{0.1139}=0.3375,\qquad \frac{\bar s}{\sqrt n}=\frac{0.3375}{3.162}=0.1067 .
\]
\textbf{گام ۴ (ضریب $t$):} درجه آزادی $n-1=9$ و از جدول $t_{9,0.025}=2.262$.

\textbf{گام ۵ (بازه):} نصف طول $2.262\times0.1067=0.2414$؛ پس
\[
\big[10.05-0.24,\ 10.05+0.24\big]=\big[9.81,\ 10.29\big].
\]
\textbf{مقایسه:} اگر (به اشتباه) $\sigma=0.3375$ را معلوم فرض می‌کردیم و از $1.96$ استفاده می‌کردیم، نصف طول $0.209$ می‌شد. بازه $t$ حدود ۱۵٪ پهن‌تر است تا عدم قطعیت در تخمین $\sigma$ را جبران کند؛ وگرنه سطح اطمینان واقعی از ۹۵٪ کمتر می‌شد.
\end{solution}
\end{example}

\begin{summarybox}[جمع‌بندی: بازه اطمینان میانگین]
\begin{center}
\renewcommand{\arraystretch}{1.5}
\begin{tabular}{c|c|c|c}
حالت & کمیت محوری & توزیع آن & بازه با سطح $1-\delta$\\\hline
$\sigma$ معلوم & $\dfrac{\bar X_n-\mu}{\sigma/\sqrt n}$ & $\mathcal N(0,1)$ & $\bar X_n\pm z_{\delta/2}\dfrac{\sigma}{\sqrt n}$\\[2mm]
$\sigma$ نامعلوم & $\dfrac{\bar X_n-\mu}{\bar S_n/\sqrt n}$ & $t_{n-1}$ & $\bar X_n\pm t_{n-1,\delta/2}\dfrac{\bar S_n}{\sqrt n}$
\end{tabular}
\end{center}
در هر دو حالت فرض گاوسی بودن داده‌ها لازم است؛ برای داده غیرگاوسی و $n$ بزرگ، به کمک قضیه حد مرکزی همین بازه‌ها تقریباً معتبرند.
\end{summarybox}

\section{تمرین‌های تکمیلی}

\begin{exercise}[بازه متقارن کوتاه‌ترین است]
در حالت گاوسی با $\sigma$ معلوم، بازه‌های نامتقارن $\big[\bar X_n-a\frac{\sigma}{\sqrt n},\ \bar X_n+b\frac{\sigma}{\sqrt n}\big]$ با $a,b>0$ را در نظر بگیرید که سطح اطمینان آن‌ها $1-\delta$ است. نشان دهید طول بازه وقتی کمینه است که $a=b$.
\end{exercise}
\begin{solution}
مانند متن، احتمال پوشش برابر $\Phi(b)-\Phi(-a)=\Phi(a)+\Phi(b)-1$ است. پس قید $\Phi(a)+\Phi(b)=2-\delta$ را داریم و می‌خواهیم $a+b$ را کمینه کنیم. از قید، $b$ تابعی از $a$ است. با مشتق‌گیری ضمنی از قید ($\phi$ چگالی گاوسی استاندارد است):
\[
\phi(a)+\phi(b)\frac{\dd b}{\dd a}=0\quad\Rightarrow\quad\frac{\dd b}{\dd a}=-\frac{\phi(a)}{\phi(b)}
\quad\Rightarrow\quad\frac{\dd(a+b)}{\dd a}=1-\frac{\phi(a)}{\phi(b)} .
\]
مشتق صفر می‌شود اگر و تنها اگر $\phi(a)=\phi(b)$، و چون $\phi$ روی $[0,\infty)$ اکیداً نزولی است، یعنی $a=b$. علامت مشتق: اگر $a<b$ آنگاه $\phi(a)>\phi(b)$ و مشتق منفی است (طول با افزایش $a$ کم می‌شود)؛ اگر $a>b$ مشتق مثبت است. پس $a=b$ کمینه است. \textbf{شهود:} بازه باید جایی باشد که چگالی بیشتر است؛ جابه‌جا کردن بازه به سمت دم، بخشی پرچگال را با بخشی کم‌چگال عوض می‌کند و برای جبران احتمال، بازه باید پهن‌تر شود.
\end{solution}

\begin{exercise}[کران اطمینان یک‌طرفه]
گاهی فقط یک کران مهم است (مثلاً «میانگین آلاینده حداکثر چقدر است؟»). در حالت گاوسی با $\sigma$ معلوم، عدد $c$ را طوری بیابید که $\Prob\big(\mu\le\bar X_n+c\frac{\sigma}{\sqrt n}\big)=1-\delta$. مقدار آن را برای $\delta=0.05$ حساب کنید.
\end{exercise}
\begin{solution}
رویداد $\mu\le\bar X_n+c\frac{\sigma}{\sqrt n}$ معادل $\frac{\bar X_n-\mu}{\sigma/\sqrt n}\ge-c$ یعنی $Z\ge-c$ است. پس
\[
\Prob(Z\ge-c)=1-\Phi(-c)=\Phi(c)=1-\delta\quad\Rightarrow\quad c=z_\delta=\Phi^{-1}(1-\delta).
\]
برای $\delta=0.05$: $c=1.645$ (کمتر از $1.96$، چون کل احتمال خطای $\delta$ در یک دم جمع شده است).
\end{solution}

\begin{exercise}[نظرسنجی: بازه اطمینان برای نسبت]
در یک نظرسنجی از $n=400$ نفر، $180$ نفر به گزینه «الف» رأی داده‌اند. (الف) با استفاده از قضیه حد مرکزی، بازه اطمینان تقریبی ۹۵٪ برای نسبت واقعی $p$ بسازید. (ب) برای اینکه نصف طول بازه ۹۵٪ (صرف‌نظر از مقدار $p$) حداکثر ۳٪ باشد، چند نفر باید پرسیده شوند؟
\end{exercise}
\begin{solution}
پاسخ هر نفر یک متغیر برنولی $X_i\sim\mathrm{Bern}(p)$ با میانگین $p$ و واریانس $p(1-p)$ است و $\hat p=\bar X_n$.

(الف) $\hat p=180/400=0.45$. طبق حد مرکزی $\frac{\hat p-p}{\sqrt{p(1-p)/n}}\approx\mathcal N(0,1)$. چون $p$ را نمی‌دانیم، در مخرج به جای آن $\hat p$ می‌گذاریم (برای $n$ بزرگ خطای ناچیزی دارد):
\[
\sqrt{\frac{\hat p(1-\hat p)}{n}}=\sqrt{\frac{0.45\times0.55}{400}}=\sqrt{0.000619}=0.0249,
\]
\[
\hat p\pm1.96\times0.0249=0.45\pm0.0488\quad\Rightarrow\quad[0.401,\ 0.499].
\]
(ب) $p(1-p)$ حداکثر در $p=\frac12$ برابر $\frac14$ است. پس با فرض بدترین حالت، باید $1.96\sqrt{\frac{1}{4n}}\le0.03$:
\[
\sqrt n\ge\frac{1.96}{2\times0.03}=32.67\quad\Rightarrow\quad n\ge1067.1\quad\Rightarrow\quad n=1068 .
\]
به همین دلیل بسیاری از نظرسنجی‌ها حدود هزار نفر نمونه دارند.
\end{solution}

\begin{exercise}[بازه اطمینان برای حد بالای توزیع یکنواخت]
نمونه‌های \lr{i.i.d.} $X_1,\ldots,X_n\sim\mathcal U(0,\theta)$ داریم و $M=\max_iX_i$ (جلسه ۲۲). (الف) نشان دهید $M/\theta$ یک کمیت محوری است. (ب) بازه اطمینانی به شکل $[M,\ cM]$ با سطح $1-\delta$ بسازید. (ج) برای $n=10$، $M=4.8$ و $\delta=0.05$ بازه را حساب کنید.
\end{exercise}
\begin{solution}
(الف) برای $0\le u\le1$: $\Prob(M/\theta\le u)=\Prob(\text{همه}\ X_i\le u\theta)=u^n$ که به $\theta$ وابسته نیست؛ پس $M/\theta$ محوری است.

(ب) چون همواره $M\le\theta$، کران پایین $M$ همیشه درست است. برای کران بالا:
\[
\Prob(\theta\le cM)=\Prob\Big(\frac M\theta\ge\frac1c\Big)=1-\Big(\frac1c\Big)^n .
\]
با قرار دادن این مقدار برابر $1-\delta$: $c^{-n}=\delta$، یعنی $c=\delta^{-1/n}$. پس بازه $\big[M,\ M\delta^{-1/n}\big]$ سطح دقیقاً $1-\delta$ دارد.

(ج) $c=0.05^{-1/10}=e^{0.1\ln20}=e^{0.2996}=1.349$؛ بازه $[4.8,\ 6.48]$. توجه کنید که بازه نامتقارن است، چون $\theta$ هرگز از $M$ کمتر نیست.
\end{solution}

\begin{exercise}[تعداد نمونه با واریانس نامعلوم]
در مثال اندازه‌گیری طول قطعه، فرض کنید می‌خواهیم نصف طول بازه ۹۵٪ حداکثر $0.1$ میلی‌متر شود. با استفاده از $\bar s=0.3375$ به عنوان حدس اولیه برای $\sigma$، تعداد نمونه لازم را تخمین بزنید.
\end{exercise}
\begin{solution}
چون $n$ احتمالاً بزرگ‌تر از ۳۰ خواهد شد، ابتدا با ضریب گاوسی تقریب می‌زنیم: $1.96\times0.3375/\sqrt n\le0.1$، پس $\sqrt n\ge6.615$ و $n\ge43.8$، یعنی $n\approx44$. با درجه آزادی $43$ داریم $t_{43,0.025}\approx2.017$ و دوباره: $\sqrt n\ge2.017\times3.375=6.81$ و $n\ge46.3$؛ پس حدود $n=47$ اندازه‌گیری لازم است. (این تنها یک تخمین است، چون $\bar S_n$ در نمونه جدید تغییر می‌کند.)
\end{solution}

\begin{exercise}[درک مفهومی]
یک دانشجو می‌گوید: «بازه اطمینان ۹۵٪ میانگین وزن، $[48.38,\ 56.22]$ است؛ پس ۹۵٪ بسته‌ها وزنی در این بازه دارند.» اشکال این جمله چیست؟
\end{exercise}
\begin{solution}
بازه اطمینان درباره \textbf{میانگین} $\mu$ است، نه درباره تک‌تک داده‌ها. وزن تک بسته‌ها انحراف معیار $\sigma=10$ دارد و حدود ۹۵٪ آن‌ها در بازه $\mu\pm1.96\times10$ (به پهنای حدود ۴۰ گرم) قرار می‌گیرند، در حالی که بازه اطمینان میانگین پهنای حدود ۸ گرم دارد، چون انحراف معیار $\bar X_n$ برابر $\sigma/\sqrt n=2$ است. افزون بر این (چنان‌که در متن دیدیم) پس از مشاهده داده، حتی گفتن «با احتمال ۹۵٪ مقدار $\mu$ در این بازه است» هم دقیق نیست؛ ۹۵٪ ویژگی روش ساخت بازه است.
\end{solution}
